\documentclass[a4paper,10pt,twocolumn]{ltjarticle}
\usepackage{kaitoushu}

\begin{document}

\problemrange{5章：三角関数 — 練習問題1・A (p.138)}

\mondai{1}{三角関数表を用いて，図の $x$，$y$ を求めよ．}
\kotae{
直角三角形なので，与えられた辺が斜辺か隣辺かを見分けて使い分ける．

\textbf{(1)} $32°$，斜辺 $3$

$x$ は $32°$ の対辺，$y$ は隣辺だから

$x=3\sin32°=3\times0.5299=1.5897\approx1.59$

$y=3\cos32°=3\times0.8480=2.5440\approx2.54$

\textbf{(2)} $57°$，隣辺 $4$

$x$ は $57°$ の対辺だから $\tan$ を使う．

$x=4\tan57°=4\times1.5399=6.1596\approx6.16$

$y$ は斜辺で，$\cos57°=\dfrac{4}{y}$ より

$y=\dfrac{4}{\cos57°}=\dfrac{4}{0.5446}=7.3448\approx7.34$
}

\mondai{2}{$\alpha$ は鈍角で $\sin\alpha=\dfrac{3}{5}$ のとき，$\dfrac{5\cos\alpha-2}{8\tan\alpha+5}$ の値を求めよ．}
\kotae{
まず $\cos\alpha$，$\tan\alpha$ を求める．

$\sin^{2}\alpha+\cos^{2}\alpha=1$ より

$\cos^{2}\alpha=1-\dfrac{9}{25}=\dfrac{16}{25}$

$\alpha$ は鈍角（$90°<\alpha<180°$）だから $\cos\alpha<0$．

$\cos\alpha=-\dfrac{4}{5}$，
$\tan\alpha=\dfrac{\sin\alpha}{\cos\alpha}=\dfrac{3/5}{-4/5}=-\dfrac{3}{4}$

これを代入して

$\dfrac{5\cos\alpha-2}{8\tan\alpha+5}
=\dfrac{5\cdot\left(-\frac{4}{5}\right)-2}{8\cdot\left(-\frac{3}{4}\right)+5}
=\dfrac{-4-2}{-6+5}=\dfrac{-6}{-1}=6$
}

\mondai{3}{次の等式を証明せよ．}
\kotae{
\textbf{(1)} $\sin^{4}\theta-\cos^{4}\theta=1-2\cos^{2}\theta$

左辺を $2$ 乗の差とみて因数分解する．

左辺 $=(\sin^{2}\theta+\cos^{2}\theta)(\sin^{2}\theta-\cos^{2}\theta)$

$\sin^{2}\theta+\cos^{2}\theta=1$ だから

$=\sin^{2}\theta-\cos^{2}\theta$

さらに $\sin^{2}\theta=1-\cos^{2}\theta$ を代入して

$=(1-\cos^{2}\theta)-\cos^{2}\theta=1-2\cos^{2}\theta=$ 右辺 \qed

\medskip
\textbf{(2)} $\tan^{2}\theta+(1-\tan^{4}\theta)\cos^{2}\theta=1$

$1+\tan^{2}\theta=\dfrac{1}{\cos^{2}\theta}$，すなわち
$(1+\tan^{2}\theta)\cos^{2}\theta=1$ を使う．

左辺の $1-\tan^{4}\theta$ を因数分解すると

$1-\tan^{4}\theta=(1-\tan^{2}\theta)(1+\tan^{2}\theta)$

よって

左辺 $=\tan^{2}\theta+(1-\tan^{2}\theta)\underbrace{(1+\tan^{2}\theta)\cos^{2}\theta}_{=1}$

$=\tan^{2}\theta+(1-\tan^{2}\theta)=1=$ 右辺 \qed
}

\mondai{4}{$\triangle$ABC について，次の問いに答えよ．}
\kotae{
\textbf{(1)} $A=85°$，$B=65°$，$c=5$ のとき，$a$，$b$ を求めよ．

まず残りの角を求める．

$C=180°-85°-65°=30°$

正弦定理 $\dfrac{a}{\sin A}=\dfrac{c}{\sin C}$ より

$a=\dfrac{c\sin A}{\sin C}=\dfrac{5\sin85°}{\sin30°}
=\dfrac{5\times0.9962}{0.5}=9.962\approx9.96$

$b=\dfrac{c\sin B}{\sin C}=\dfrac{5\sin65°}{\sin30°}
=\dfrac{5\times0.9063}{0.5}=9.063\approx9.06$

\medskip
\textbf{(2)} $a=4$，$b=3$，$c=\sqrt{37}$ のとき，$C$ を求めよ．

3辺がわかっているので余弦定理を使う．

$\cos C=\dfrac{a^{2}+b^{2}-c^{2}}{2ab}
=\dfrac{16+9-37}{2\cdot4\cdot3}=\dfrac{-12}{24}=-\dfrac{1}{2}$

$0°<C<180°$ より

$\therefore C=120°$
}

\mondai{5}{$\triangle$ABC において，等式 $\cos C=\dfrac{\sin A}{\sin B}$ が成り立つとき，次の問いに答えよ．ただし，$R$ は $\triangle$ABC の外接円の半径とする．}
\kotae{
\textbf{(1)} 等式を $a$，$b$，$c$，$R$ の関係式に直せ．

正弦定理 $\dfrac{a}{\sin A}=\dfrac{b}{\sin B}=2R$ より

$\sin A=\dfrac{a}{2R}$，$\sin B=\dfrac{b}{2R}$

よって右辺は

$\dfrac{\sin A}{\sin B}=\dfrac{a/(2R)}{b/(2R)}=\dfrac{a}{b}$

（$2R$ が約分され，$R$ は消える．）

左辺は余弦定理より

$\cos C=\dfrac{a^{2}+b^{2}-c^{2}}{2ab}$

したがって

$\dfrac{a^{2}+b^{2}-c^{2}}{2ab}=\dfrac{a}{b}$

両辺に $2ab$ を掛けて

$a^{2}+b^{2}-c^{2}=2a^{2}$

$\therefore b^{2}=a^{2}+c^{2}$

\medskip
\textbf{(2)} $\triangle$ABC はどんな三角形か．

(1) より $b^{2}=a^{2}+c^{2}$．
三平方の定理の逆から，最大辺 $b$ に向かい合う角 $B$ が直角である．

$\therefore \angle B=90°$ の直角三角形
}

\mondai{6}{図の台形について，辺 AB の長さおよび面積 $S$ を求めよ．
（$AD\parallel BC$，$AD=4$，$BC=8$，$\angle B=70°$，$\angle C=57°$）}
\kotae{
A，D から辺 BC に垂線 AH，DK を下ろす．
$AD\parallel BC$ だから $AH=DK$ で，これを高さ $h$ とおく．

$HK=AD=4$ であり，直角三角形 ABH，DKC において

$BH=\dfrac{h}{\tan70°}$，$KC=\dfrac{h}{\tan57°}$

$BC=BH+HK+KC$ だから

$\dfrac{h}{\tan70°}+4+\dfrac{h}{\tan57°}=8$

$h\left(\dfrac{1}{\tan70°}+\dfrac{1}{\tan57°}\right)=4$

三角関数表より $\tan70°=2.7475$，$\tan57°=1.5399$ なので

$\dfrac{1}{2.7475}+\dfrac{1}{1.5399}=0.3640+0.6494=1.0134$

$h=\dfrac{4}{1.0134}=3.9471\approx3.95$

次に AB を求める．直角三角形 ABH で
$\sin70°=\dfrac{h}{AB}$ だから

$AB=\dfrac{h}{\sin70°}=\dfrac{3.9471}{0.9397}=4.2004\approx4.20$

面積は台形の公式より

$S=\dfrac{1}{2}(AD+BC)\times h
=\dfrac{1}{2}(4+8)\times3.9471=6\times3.9471=23.68$

$\therefore AB\approx4.20$，$S\approx23.7$
}

\end{document}
